\begin{center}
    \textbf{Resumo}
\end{center}

\noindent O uso de estruturas de concreto protendido cresceu e ganhou notoriedade após a Segunda Guerra Mundial devido ao desenvolvimento de aços de alta resistência, dado sua versatilidade construtiva para estruturas que exigem maiores vãos livres. Dentre os cálculos feitos no dimensionamento de projetos desse método construtivo estão as perdas de protensão, que se mostram um verdadeiro desafio dado sua extensão e complexidade. Embora o cálculo preciso dessas perdas seja crucial para a otimização estrutural, ele é frequentemente negligenciado, resultando em superdimensionamento e aumento de custos. Ferramentas computacionais existentes, muitas vezes opacas e pouco explicáveis, não abordam adequadamente esse aspecto, levando ao uso de constantes simplificadas. Dado esse cenário, o presente estudo propõe o desenvolvimento de uma ferramenta computacional acessível, gratuita e de código aberto, escrita na linguagem de programação \textit{Python}, para facilitar o cálculo transparente e modificado de perdas de protensão, promovendo a clareza, eficiência e economia no dimensionamento de estruturas de concreto protendido.
\vspace{0.2cm}

\noindent\textbf{Palavras-chave}: concreto protendido; perda de protensão; ferramenta computacional.